\begin{frame}[fragile]{실습: GBDT vs 랜덤 포레스트 (같은 데이터)}
  7.2절의 유방암 데이터:
\begin{lstlisting}
from sklearn.ensemble import RandomForestClassifier, GradientBoostingClassifier

rf = RandomForestClassifier(n_estimators=100, random_state=0).fit(X_train, y_train)
gbdt = GradientBoostingClassifier(n_estimators=100, learning_rate=0.1,
                                  random_state=0).fit(X_train, y_train)
print("RandomForest:", round(rf.score(X_test, y_test), 3))
print("GBDT:        ", round(gbdt.score(X_test, y_test), 3))
\end{lstlisting}
  \vskip0.3em
  {\footnotesize \texttt{RandomForest: 0.959 \quad GBDT: 0.977} \quad
  상위 3 중요도: \texttt{worst concave points 0.445},
  \texttt{mean concave points 0.275}, \texttt{worst perimeter 0.071}}
  \vskip0.4em
  \begin{itemize}
    \item GBDT가 RF보다 \textbf{근소하게 더 높음}(0.977 vs 0.959) —
      ``독립적으로 합치는 것''보다 ``\textbf{잔차를 순차적으로 쫓는
      것}''이 조금 더 정교하게 맞춤
    \item \texttt{feature\_importances\_}는 \textbf{모델 전체}의 설명이지
      ``이 환자 한 명의 예측'' 설명은 아님 $\to$ 개별 예측 하나를
      설명하려면 \textbf{SHAP}이 필요
  \end{itemize}
\end{frame}
